\section{Cloud Begriffe}
\subsection{Cloud consumer}
Ein Cloud Consumer ist der Nutzer von durch eine Cloud bereitgestellten Diensten. Dieser entscheidet über seine eigenen Abonnements. Ein Beispiel ist wenn man sich für ein Claude – Also den KI Chatbot – Abo entscheidet ist man ein Cloud Consumer. Dasselbe gilt für iCloud, Spotify, Apple Music und Netflix. Im Enterprise Kontext ist das noch mal anders dort sind typische Abonnements Microsoft 365, Salesforce oder auch Abonnements bei Cloud Providern wie bei Azure.
\subsection{Cloud Provider}
Ein Cloud Provider stellt die Infrastruktur und die Services einer Cloud bereit. Abseits davon sind die Provider für Betrieb und Sicherheit der Dienste verantwortlich. Beispiele sind die Hyperscaler wie Azure, AWS und Google aber auch europäische Anbieter wie beispielsweise Hetzner sind Cloud Provider.
\subsection{Cloudauditor}
Ein Cloud Auditor prüft Rechenzentren von Cloud Providern und ob diese sich an den Datenschutz und Compliance halten.
\subsection{Cloud broker}
Ein Cloud Broker ist ein unabhängiger Anbieter der zwischen Cloud Consumer und Provider kommuniziert. Ein Beispiel wäre ein Systemhaus das Lizenzen einer Cloud kauft und Consumer diese samt Support und Konfiguration weiter verkauft.
\subsection{Cloud carrier}
Ein Cloud Carrier ist für eine sichere und stabile Datenübertragung von Cloud Consumer zu Cloud Provider verantwortlich.
